\subsection{Ampliación de Conjuntos de Datos de Contexto por Ingestión de Documentos}

En el punto 4.1 se crearon conjuntos de datos específicos para la problemática del problema. El paso siguiente es extender los experimentos hacia un conjunto de datos más extenso y con una estructura extraída de un origen de datos real.

En este paso se seleccionaron resúmenes de tarjetas de crédito para crear un conjunto de datos de consumos sobre el que continuar el análisis.

El procedimiento consistía en los siguientes pasos:

\begin{enumerate}
    \item Transformar los archivos PDF a texto plano.
    \item Interpretar el texto y extraer entidades.
    \item Ordenar las entidades para guardarlas en una base vectorial.
    \item Testear la correcta recuperación del dato mediante lenguaje natural por el método de \textit{Retrieval Augmented Generation} (RAG) simple, esto es, sin reranking posterior.
\end{enumerate}

Al momento de desarrollar esta sección del trabajo los modelos multimodales, aquellos que pueden procesar imágenes y texto conjuntamente, no existían y las técnicas necesarias para extraer correctamente información sobre consumos de un resumen de tarjeta de crédito requerían de un esfuerzo comparable con todo el objeto de estudio de esta tesis.

Los intentos fallidos de extracción no están publicados ya que no aportan a la descripción del avance del desarrollo.

En este punto junto al cuerpo docente se decidió tomar un camino diferente, centrado en crear toda la información de una misma persona de manera sintética, con un modelo de lenguaje de primer nivel como GPT-4, el líder en ese momento. 